\chapter{Introdução}\label{cap:introducao}

Enquanto modalidade com a maior projeção mediática e económica a nível mundial, o futebol tem
assistido a uma adoção acelerada de soluções de visão por computador e aprendizagem automática.
Clubes, federações e entidades de \emph{broadcast} apostam progressivamente em ferramentas que
permitam modelar o comportamento das equipas, antever dinâmicas de jogo e sustentar a decisão
tática em evidência quantitativa, em detrimento de uma análise estritamente empírica. Neste
contexto, a extração automática de informação estruturada a partir do sinal visual de transmissão
assumiu-se como um fator claramente diferenciador.

É neste contexto que se enquadra a reconstrução do estado de jogo (\emph{Game State
Reconstruction}, GSR), formalizada por \citet{somers2024soccernet}. Trata-se de uma tarefa de
natureza integradora, cujo objetivo é reconstruir, a partir de um vídeo de uma única câmara, uma
representação bidimensional do campo onde cada interveniente é caracterizado pela sua posição, pela
equipa a que pertence, pela função que desempenha (jogador, guarda-redes ou árbitro) e pelo número
de camisola. Reunida num único mapa, esta informação permite seguir cada jogador ao longo do jogo,
ler as formações e analisar as decisões táticas.

A reconstrução do estado de jogo constitui uma tarefa intrinsecamente complexa, uma vez que articula
num único \emph{pipeline} múltiplas sub-tarefas de elevada exigência individual. Inclui a deteção e
o \emph{tracking} dos jogadores, a sua reidentificação após oclusões ou saídas de campo, a
calibração da câmara para o mapeamento no relvado e, por fim, a sua identificação individual. É
precisamente nesta última etapa que se concentram dois dos componentes mais vulneráveis do
\emph{pipeline}, o reconhecimento do número de camisola e a atribuição da equipa. Em ambos os
casos, o sinal visual disponível é frequentemente escasso e ruidoso, caracterizando-se por números
de pequena dimensão, oclusões parciais, ângulos desfavoráveis e equipamentos cujas cores nem sempre
oferecem um contraste claro. Ainda assim, é destes dois módulos que depende a capacidade de associar uma identidade única a cada
interveniente, e sem eles todo o processamento anterior perde grande parte do seu valor analítico.

Apesar da sua importância, a identificação dos jogadores tem recebido uma atenção substancialmente
menor do que o \emph{tracking}, que concentra a maior parte da investigação na análise e na recolha
de dados estatísticos do jogo de futebol. Esta assimetria acarreta limitações claras. Por mais
robusto que seja, um sistema capaz de seguir os intervenientes no relvado mas incapaz de os
identificar falha em responder à questão mais elementar de qualquer análise desportiva, a de saber
\emph{quem} é cada jogador em campo. É precisamente esta lacuna que a presente dissertação se propõe
colmatar.

O trabalho centra-se, então, nos dois componentes de identificação que pior se comportam na
solução de referência do desafio SoccerNet, a leitura do número de camisola e a atribuição de cada
jogador à sua equipa e ao respetivo lado do campo. Para os melhorar, alteram-se quatro dos nove
módulos do \emph{pipeline}. O leitor de números é substituído, os módulos da equipa e do lado do
campo são reescritos e a reidentificação ganha um descritor da cor do equipamento, sem mudar o que
já produzia. A deteção, o \emph{tracking}, a calibração e a classificação da função ficam como estavam,
para que qualquer diferença de desempenho possa ser atribuída, sem margem para dúvidas, ao que de
facto mudou. Na leitura do
número, o reconhecimento ótico de caracteres dá lugar a um modelo multimodal de visão e linguagem,
que lê apenas a zona das costas do jogador. As leituras de cada trajetória são depois combinadas,
mas só contam aquelas em que o próprio modelo diz estar seguro. Na atribuição da equipa, corrigem-se erros no agrupamento por
aparência e no modo como cada grupo é associado a um lado do campo, e afinam-se os limiares de
decisão em conjunto sobre os dados de teste e de validação, com a regra de nunca piorar um caso que
já estava certo.

Avaliadas nos $159$ vídeos anotados do conjunto de dados, contra a solução de referência executada
nas mesmas condições e sobre os mesmos \emph{tracklets}, estas alterações traduzem-se em ganhos
claros, tanto nas duas tarefas isoladamente como na métrica oficial do desafio, o GS-HOTA, que é
ainda decomposta atributo a atributo para que o contributo de cada módulo possa ser lido
separadamente. O conjunto de treino, que não interveio em nenhuma decisão de conceção, produz
ganhos da mesma ordem dos do conjunto de teste, o que afasta a hipótese de os resultados decorrerem
de um sobreajuste aos dados sobre os quais o sistema foi afinado. E tão importante como os ganhos
foi perceber até onde eles vão. Um conjunto de experiências controladas mostrou que, na leitura do
número, o desempenho que ainda falta está limitado pela informação que a própria imagem contém, o
que remete o próximo avanço para etapas anteriores do processamento. Na atribuição da equipa, a
inspeção visual de todos os vídeos caracterizou as falhas que subsistem, incluindo uma regressão
face à referência com causa identificada, e revelou de passagem um defeito de calibração na
configuração distribuída da própria solução de referência.

\section{Estrutura do Documento}\label{sec:estrutura}

A dissertação está organizada em cinco capítulos. O presente capítulo enquadra o problema e delimita
o âmbito do trabalho. O Capítulo~\ref{cap:estado-arte} revê o estado da arte sobre reconhecimento de
número de camisola e atribuição da equipa, e percorre os métodos apresentados às sucessivas edições
do desafio SoccerNet. O Capítulo~\ref{cap:metodologia} descreve a metodologia proposta e a sua
implementação. O Capítulo~\ref{cap:resultados} reporta e discute os resultados experimentais. Por
fim, o Capítulo~\ref{cap:conclusao} sintetiza os principais contributos e aponta direções de
trabalho futuro.
